\chapter*{Future Scope}

While the current version of ProblySUS establishes a strong multi-signal foundation for URL scam detection, several strategic enhancements are planned to extend its reach and capability.

\begin{itemize}

\item \textbf{Machine Learning-Based Scoring:}  
The current rule-based scoring engine could be augmented or replaced with a trained machine learning classifier such as Random Forest or XGBoost. Training the model on a labelled dataset of phishing and legitimate URLs could improve detection accuracy beyond fixed threshold penalties.

\item \textbf{Public Cloud Deployment:}  
Deploying the backend to a cloud platform such as AWS, Google Cloud Platform, or Render would make ProblySUS publicly accessible without requiring users to run a local server. This would significantly widen its user base and allow the browser extension to be used by anyone.

\item \textbf{Real-Time Passive Protection via Extension:}  
The browser extension currently requires users to manually click the ``Scan'' button. A future version could monitor tab navigation events and proactively warn users before they land on a risky page, providing real-time safe-browsing protection.

\item \textbf{Expanded Threat Intelligence Feeds:}  
Integrating additional threat intelligence sources such as PhishTank, OpenPhish, or Emerging Threats would improve blacklist coverage, particularly for newly launched phishing campaigns that may not yet appear in the URLhaus dataset.

\item \textbf{Historical Domain Reputation Analysis:}  
Cross-referencing domains against historical reputation databases such as the Wayback Machine and historical hosting records could provide additional trust signals beyond the current WHOIS-based domain age analysis.

\item \textbf{API Rate Limiting and Authentication:}  
For public deployment, implementing API key authentication and per-key rate limiting would prevent abuse of the service and ensure fair resource allocation among users.

\item \textbf{User Reporting and Feedback Loop:}  
Introducing a user feedback system would allow users to report false positives and false negatives. This feedback could be used to continuously refine scoring thresholds and improve detection accuracy.

\item \textbf{Multilingual URL and Content Analysis:}  
Extending the content analysis module to detect phishing indicators across multiple languages would enhance detection accuracy for non-English phishing campaigns, which are increasingly common.

\end{itemize}

Robust privacy controls will also be prioritized as the system evolves. These include enforcing HTTPS for all public API endpoints, encrypting stored result logs, and ensuring compliance with modern data protection standards when transitioning from local deployment to a cloud-based infrastructure.